% Source: https://artofproblemsolving.com/community/c3h3589608_variety_of_geometry
% Stable identity: aops-topic-3589608-q22
Scale and orient the figure so $D=(0,0)$, $A=(1,0)$, and $B=(0,\sqrt3)$; this follows from the right triangle $ADB$ and $\angle BAD=60^\circ$. Convexity and $\angle ADC=135^\circ$ give $C=(-c,c)$ with $c>0$. In triangle $BCD$, the angles at $D,C,B$ are $45^\circ,60^\circ,75^\circ$. The sine rule gives $\dfrac{c\sqrt2}{\sin75^\circ}=\dfrac{\sqrt3}{\sin60^\circ}=2$, so $c=\dfrac{\sqrt3+1}{2}$. The line from $(1,0)$ to $(-c,c)$ meets $x=0$ at height $\dfrac{c}{1+c}=\dfrac{1}{\sqrt3}$. Thus $DO=\dfrac{1}{\sqrt3}$ and $OB=\dfrac{2}{\sqrt3}$, and the ratio is $\boxed{\dfrac{1}{2}}$.
\begin{asy}
import geometry;
size(240);
real c=(sqrt(3)+1)/2;pair D=(0,0),A=(1,0),B=(0,sqrt(3)),C=(-c,c),O=(0,1/sqrt(3));draw(A--B--C--D--cycle);draw(A--C^^B--D,dashed);label("$A$",A,E);label("$B$",B,N);label("$C$",C,W);label("$D$",D,S);label("$O$",O,NE);
\end{asy}
